\section{Geometric Orbit and Stabilizer Tables}
\label{sec:geometric-stabilizers}

For this section let \(k=\overline{\mathbf F}_q\) and \(V=k^n\) with
\(n=4,5,6,7\). The following theorem records the geometric orbit
representatives together with the exact kernel and quotient data attached to
their stabilizers.

\begin{theorem}
\label{thm:geometric-tables}
For \(n=4,5,6,7\), the \(\GL(V)\)-orbits on \(\Gr(2,\wedge^2 V)\) are exactly
the rows listed in the geometric orbit tables below. For each representative
\(L\), the stabilizer is determined by the exact sequence
\[
  1\to K_L\to \Stab_{\GL(V)}(L)\xrightarrow{\rho_L} Q_L\to 1,
\]
with \(K_L\) and \(Q_L\) given by the corresponding row of the geometric
stabilizer tables.
\end{theorem}

\begin{proof}
The kernel and quotient statements are verified row by row in the companion
Lean~4 formalization; see Subsection~\ref{subsec:formalization-protocol} and
Appendix~\ref{app:lean-crosswalk}. The completeness of the geometric list is
deduced later in Corollary~\ref{cor:geometric-list-complete-rewrite} from the
finite-field orbit-count identity.
\end{proof}

In the quotient column, \(B\), \(T\), and \(N_T\) denote the full preimages in
\(\GL(L)\cong \GL_2(k)\) of a Borel subgroup, a maximal torus, and the
normalizer of a maximal torus in \(\PGL(L)\). The entry ``\(S_3\) preimage''
denotes the full preimage of \(S_3\le \PGL(L)\).

\subsection{Geometric orbit representatives}

\begingroup
\small
% This file is generated by magma/examples/generate_geometric_orbit_tables_tex.m.

\paragraph{$n=4$.}
\begin{longtable}{@{}c p{0.26\textwidth} p{0.25\textwidth} p{0.25\textwidth}@{}}
\toprule
No. & Type & $\omega_1$ & $\omega_2$ \\
\midrule
\endfirsthead
\toprule
No. & Type & $\omega_1$ & $\omega_2$ \\
\midrule
\endhead
\midrule
\multicolumn{4}{r}{Continued on next page}\\
\midrule
\endfoot
\bottomrule
\endlastfoot
1 & \(J_{a,2}\) & $e_{1} \wedge e_{3} + e_{2} \wedge e_{4}$ & $e_{1} \wedge e_{4}$ \\
2 & \(J_{a,1} + J_{b,1}\) & $e_{1} \wedge e_{2} + e_{3} \wedge e_{4}$ & $e_{3} \wedge e_{4}$ \\
3 & \(S_1 + S_2\) & $e_{2} \wedge e_{4}$ & $e_{3} \wedge e_{4}$ \\
\end{longtable}

\paragraph{$n=5$.}
\begin{longtable}{@{}c p{0.26\textwidth} p{0.25\textwidth} p{0.25\textwidth}@{}}
\toprule
No. & Type & $\omega_1$ & $\omega_2$ \\
\midrule
\endfirsthead
\toprule
No. & Type & $\omega_1$ & $\omega_2$ \\
\midrule
\endhead
\midrule
\multicolumn{4}{r}{Continued on next page}\\
\midrule
\endfoot
\bottomrule
\endlastfoot
1 & \(S_3\) & $e_{1} \wedge e_{4} + e_{2} \wedge e_{5}$ & $e_{2} \wedge e_{4} + e_{3} \wedge e_{5}$ \\
2 & \(S_2 + S_1^2\) & $e_{1} \wedge e_{3}$ & $e_{2} \wedge e_{3}$ \\
3 & \(S_2 + J_{a,1}\) & $e_{1} \wedge e_{3} + e_{4} \wedge e_{5}$ & $e_{2} \wedge e_{3}$ \\
4 & \(S_1 + J_{a,2}\) & $e_{2} \wedge e_{4} + e_{3} \wedge e_{5}$ & $e_{2} \wedge e_{5}$ \\
5 & \(S_1 + J_{a,1} + J_{b,1}\) & $e_{2} \wedge e_{3} + e_{4} \wedge e_{5}$ & $e_{4} \wedge e_{5}$ \\
\end{longtable}

\paragraph{$n=6$.}
\begin{longtable}{@{}c p{0.26\textwidth} p{0.25\textwidth} p{0.25\textwidth}@{}}
\toprule
No. & Type & $\omega_1$ & $\omega_2$ \\
\midrule
\endfirsthead
\toprule
No. & Type & $\omega_1$ & $\omega_2$ \\
\midrule
\endhead
\midrule
\multicolumn{4}{r}{Continued on next page}\\
\midrule
\endfoot
\bottomrule
\endlastfoot
1 & \(J_{a,3}\) & $e_{1} \wedge e_{4} + e_{2} \wedge e_{5} + e_{3} \wedge e_{6}$ & $e_{1} \wedge e_{5} + e_{2} \wedge e_{6}$ \\
2 & \(J_{a,2} + J_{a,1}\) & $e_{1} \wedge e_{3} + e_{2} \wedge e_{4} + e_{5} \wedge e_{6}$ & $e_{1} \wedge e_{4}$ \\
3 & \(J_{a,2} + J_{b,1}\) & $e_{1} \wedge e_{3} + e_{2} \wedge e_{4} + e_{5} \wedge e_{6}$ & $e_{1} \wedge e_{4} + e_{5} \wedge e_{6}$ \\
4 & \(2J_{a,1} + J_{b,1}\) & $e_{1} \wedge e_{2} + e_{3} \wedge e_{4} + e_{5} \wedge e_{6}$ & $e_{5} \wedge e_{6}$ \\
5 & \(J_{a,1} + J_{b,1} + J_{c,1}\) & $e_{1} \wedge e_{2} + e_{3} \wedge e_{4}$ & $e_{3} \wedge e_{4} + e_{5} \wedge e_{6}$ \\
6 & \(S_1^2 + J_{a,2}\) & $e_{3} \wedge e_{5} + e_{4} \wedge e_{6}$ & $e_{3} \wedge e_{6}$ \\
7 & \(S_1^2 + J_{a,1} + J_{b,1}\) & $e_{3} \wedge e_{4} + e_{5} \wedge e_{6}$ & $e_{5} \wedge e_{6}$ \\
8 & \(S_1 + S_2 + J_{a,1}\) & $e_{2} \wedge e_{4} + e_{5} \wedge e_{6}$ & $e_{3} \wedge e_{4}$ \\
9 & \(S_1^3 + S_2\) & $e_{4} \wedge e_{6}$ & $e_{5} \wedge e_{6}$ \\
10 & \(S_1 + S_3\) & $e_{2} \wedge e_{5} + e_{3} \wedge e_{6}$ & $e_{3} \wedge e_{5} + e_{4} \wedge e_{6}$ \\
11 & \(S_2^2\) & $e_{1} \wedge e_{3} + e_{4} \wedge e_{6}$ & $e_{2} \wedge e_{3} + e_{5} \wedge e_{6}$ \\
\end{longtable}

\paragraph{$n=7$.}
\begin{longtable}{@{}c p{0.26\textwidth} p{0.25\textwidth} p{0.25\textwidth}@{}}
\toprule
No. & Type & $\omega_1$ & $\omega_2$ \\
\midrule
\endfirsthead
\toprule
No. & Type & $\omega_1$ & $\omega_2$ \\
\midrule
\endhead
\midrule
\multicolumn{4}{r}{Continued on next page}\\
\midrule
\endfoot
\bottomrule
\endlastfoot
1 & \(S_4\) & $e_{1} \wedge e_{5} + e_{2} \wedge e_{6} + e_{3} \wedge e_{7}$ & $e_{2} \wedge e_{5} + e_{3} \wedge e_{6} + e_{4} \wedge e_{7}$ \\
2 & \(S_3 + S_1^2\) & $e_{1} \wedge e_{4} + e_{2} \wedge e_{5}$ & $e_{2} \wedge e_{4} + e_{3} \wedge e_{5}$ \\
3 & \(S_2^2 + S_1\) & $e_{1} \wedge e_{3} + e_{4} \wedge e_{6}$ & $e_{2} \wedge e_{3} + e_{5} \wedge e_{6}$ \\
4 & \(S_2 + S_1^4\) & $e_{1} \wedge e_{3}$ & $e_{2} \wedge e_{3}$ \\
5 & \(S_3 + J_{a,1}\) & $e_{1} \wedge e_{4} + e_{2} \wedge e_{5} + e_{6} \wedge e_{7}$ & $e_{2} \wedge e_{4} + e_{3} \wedge e_{5}$ \\
6 & \(S_2 + S_1^2 + J_{a,1}\) & $e_{1} \wedge e_{3} + e_{6} \wedge e_{7}$ & $e_{2} \wedge e_{3}$ \\
7 & \(S_2 + J_{a,2}\) & $e_{1} \wedge e_{3} + e_{4} \wedge e_{6} + e_{5} \wedge e_{7}$ & $e_{2} \wedge e_{3} + e_{4} \wedge e_{7}$ \\
8 & \(S_2 + 2J_{a,1}\) & $e_{1} \wedge e_{3} + e_{4} \wedge e_{5} + e_{6} \wedge e_{7}$ & $e_{2} \wedge e_{3}$ \\
9 & \(S_2 + J_{a,1} + J_{b,1}\) & $e_{1} \wedge e_{3} + e_{4} \wedge e_{5} + e_{6} \wedge e_{7}$ & $e_{2} \wedge e_{3} + e_{6} \wedge e_{7}$ \\
10 & \(S_1^3 + J_{a,2}\) & $e_{4} \wedge e_{6} + e_{5} \wedge e_{7}$ & $e_{4} \wedge e_{7}$ \\
11 & \(S_1^3 + J_{a,1} + J_{b,1}\) & $e_{4} \wedge e_{5} + e_{6} \wedge e_{7}$ & $e_{6} \wedge e_{7}$ \\
12 & \(S_1 + J_{a,3}\) & $e_{2} \wedge e_{5} + e_{3} \wedge e_{6} + e_{4} \wedge e_{7}$ & $e_{2} \wedge e_{6} + e_{3} \wedge e_{7}$ \\
13 & \(S_1 + J_{a,2} + J_{a,1}\) & $e_{2} \wedge e_{4} + e_{3} \wedge e_{5} + e_{6} \wedge e_{7}$ & $e_{2} \wedge e_{5}$ \\
14 & \(S_1 + J_{a,2} + J_{b,1}\) & $e_{2} \wedge e_{4} + e_{3} \wedge e_{5} + e_{6} \wedge e_{7}$ & $e_{2} \wedge e_{5} + e_{6} \wedge e_{7}$ \\
15 & \(S_1 + 2J_{a,1} + J_{b,1}\) & $e_{2} \wedge e_{3} + e_{4} \wedge e_{5} + e_{6} \wedge e_{7}$ & $e_{6} \wedge e_{7}$ \\
16 & \(S_1 + J_{a,1} + J_{b,1} + J_{c,1}\) & $e_{2} \wedge e_{3} + e_{4} \wedge e_{5}$ & $e_{4} \wedge e_{5} + e_{6} \wedge e_{7}$ \\
\end{longtable}

\endgroup

\subsection{Geometric stabilizers}

\begingroup
\small
% This file is generated by magma/examples/generate_geometric_stabilizer_tables_tex.m.

\paragraph{$n=4$.}
\begin{longtable}{@{}c p{0.28\textwidth} p{0.47\textwidth} p{0.12\textwidth}@{}}
\toprule
No. & Type & $K_L$ in $\Stab_{\GL(V)}(L)=K_L\rtimes Q_L$ & $Q_L$ \\
\midrule
\endfirsthead
\toprule
No. & Type & $K_L$ in $\Stab_{\GL(V)}(L)=K_L\rtimes Q_L$ & $Q_L$ \\
\midrule
\endhead
\midrule
\multicolumn{4}{r}{Continued on next page}\\
\midrule
\endfoot
\bottomrule
\endlastfoot
1 & \(J_{a,2}\) & $U_3\rtimes \SL_2(k)$ & $B$ \\
2 & \(J_{a,1} + J_{b,1}\) & $\SL_2(k)\times \SL_2(k)$ & $N_T$ \\
3 & \(S_1 + S_2\) & $U_5\rtimes (\Gm\times \Gm)$ & $\GL_2(k)$ \\
\end{longtable}

\paragraph{$n=5$.}
\begin{longtable}{@{}c p{0.28\textwidth} p{0.47\textwidth} p{0.12\textwidth}@{}}
\toprule
No. & Type & $K_L$ in $\Stab_{\GL(V)}(L)=K_L\rtimes Q_L$ & $Q_L$ \\
\midrule
\endfirsthead
\toprule
No. & Type & $K_L$ in $\Stab_{\GL(V)}(L)=K_L\rtimes Q_L$ & $Q_L$ \\
\midrule
\endhead
\midrule
\multicolumn{4}{r}{Continued on next page}\\
\midrule
\endfoot
\bottomrule
\endlastfoot
1 & \(S_3\) & $U_4\rtimes \Gm$ & $\GL_2(k)$ \\
2 & \(S_2 + S_1^2\) & $U_8\rtimes (\GL_2(k)\times \Gm)$ & $\GL_2(k)$ \\
3 & \(S_2 + J_{a,1}\) & $U_4\rtimes (\Gm\times \SL_2(k))$ & $B$ \\
4 & \(S_1 + J_{a,2}\) & $U_7\rtimes (\Gm\times \SL_2(k))$ & $B$ \\
5 & \(S_1 + J_{a,1} + J_{b,1}\) & $U_4\rtimes (\Gm\times \SL_2(k)\times \SL_2(k))$ & $N_T$ \\
\end{longtable}

\paragraph{$n=6$.}
\begin{longtable}{@{}c p{0.28\textwidth} p{0.47\textwidth} p{0.12\textwidth}@{}}
\toprule
No. & Type & $K_L$ in $\Stab_{\GL(V)}(L)=K_L\rtimes Q_L$ & $Q_L$ \\
\midrule
\endfirsthead
\toprule
No. & Type & $K_L$ in $\Stab_{\GL(V)}(L)=K_L\rtimes Q_L$ & $Q_L$ \\
\midrule
\endhead
\midrule
\multicolumn{4}{r}{Continued on next page}\\
\midrule
\endfoot
\bottomrule
\endlastfoot
1 & \(J_{a,3}\) & $U_6\rtimes \SL_2(k)$ & $B$ \\
2 & \(J_{a,2} + J_{a,1}\) & $U_7\rtimes (\SL_2(k)\times \SL_2(k))$ & $B$ \\
3 & \(J_{a,2} + J_{b,1}\) & $U_3\rtimes (\SL_2(k)\times \SL_2(k))$ & $T$ \\
4 & \(2J_{a,1} + J_{b,1}\) & $\Sp_4(k)\times \SL_2(k)$ & $T$ \\
5 & \(J_{a,1} + J_{b,1} + J_{c,1}\) & $\SL_2(k)\times \SL_2(k)\times \SL_2(k)$ & $S_3$ preimage \\
6 & \(S_1^2 + J_{a,2}\) & $U_{11}\rtimes (\GL_2(k)\times \SL_2(k))$ & $B$ \\
7 & \(S_1^2 + J_{a,1} + J_{b,1}\) & $U_8\rtimes (\GL_2(k)\times \SL_2(k)\times \SL_2(k))$ & $N_T$ \\
8 & \(S_1 + S_2 + J_{a,1}\) & $U_9\rtimes (\Gm\times \Gm\times \SL_2(k))$ & $B$ \\
9 & \(S_1^3 + S_2\) & $U_{11}\rtimes (\GL_3(k)\times \Gm)$ & $\GL_2(k)$ \\
10 & \(S_1 + S_3\) & $U_8\rtimes (\Gm\times \Gm)$ & $\GL_2(k)$ \\
11 & \(S_2^2\) & $U_6\rtimes \GL_2(k)$ & $\GL_2(k)$ \\
\end{longtable}

\paragraph{$n=7$.}
\begin{longtable}{@{}c p{0.28\textwidth} p{0.47\textwidth} p{0.12\textwidth}@{}}
\toprule
No. & Type & $K_L$ in $\Stab_{\GL(V)}(L)=K_L\rtimes Q_L$ & $Q_L$ \\
\midrule
\endfirsthead
\toprule
No. & Type & $K_L$ in $\Stab_{\GL(V)}(L)=K_L\rtimes Q_L$ & $Q_L$ \\
\midrule
\endhead
\midrule
\multicolumn{4}{r}{Continued on next page}\\
\midrule
\endfoot
\bottomrule
\endlastfoot
1 & \(S_4\) & $\Ga^{14}\rtimes (\Gm\times \Gm)$ & $\GL_2(k)$ \\
2 & \(S_3 + S_1^2\) & $\Ga^{10}\rtimes (\Gm\times (U_4\rtimes \Gm))$ & $\GL_2(k)$ \\
3 & \(S_2^2 + S_1\) & $\Ga^6\rtimes (\Gm\times (U_6\rtimes \GL_2(k)))$ & $\GL_2(k)$ \\
4 & \(S_2 + S_1^4\) & $\Ga^{12}\rtimes (\GL_4(k)\times (\Ga^2\rtimes \Gm))$ & $\GL_2(k)$ \\
5 & \(S_3 + J_{a,1}\) & $\Ga^{10}\rtimes (\SL_2(k)\times (U_4\rtimes \Gm))$ & $B$ \\
6 & \(S_2 + S_1^2 + J_{a,1}\) & $\Ga^{10}\rtimes (\Gm\times (U_4\rtimes (\Gm\times \SL_2(k))))$ & $B$ \\
7 & \(S_2 + J_{a,2}\) & $\Ga^6\rtimes (\Gm\times (U_3\rtimes \SL_2(k)))$ & $B$ \\
8 & \(S_2 + 2J_{a,1}\) & $\Ga^6\rtimes (\Gm\times \Sp_4(k))$ & $B$ \\
9 & \(S_2 + J_{a,1} + J_{b,1}\) & $\Ga^6\rtimes (\Gm\times (\SL_2(k)\times \SL_2(k)))$ & $N_T$ \\
10 & \(S_1^3 + J_{a,2}\) & $\Ga^{12}\rtimes (\GL_3(k)\times \SL_2(k[t]/(t^2)))$ & $B$ \\
11 & \(S_1^3 + J_{a,1} + J_{b,1}\) & $\Ga^{12}\rtimes (\GL_3(k)\times (\SL_2(k)\times \SL_2(k)))$ & $N_T$ \\
12 & \(S_1 + J_{a,3}\) & $\Ga^6\rtimes (\Gm\times \SL_2(k[t]/(t^3)))$ & $B$ \\
13 & \(S_1 + J_{a,2} + J_{a,1}\) & $\Ga^6\rtimes (\Gm\times (U_{2,1}(k)\rtimes (\SL_2(k)\times \SL_2(k))))$ & $B$ \\
14 & \(S_1 + J_{a,2} + J_{b,1}\) & $\Ga^6\rtimes (\Gm\times (\SL_2(k)\times \SL_2(k)))$ & $T$ \\
15 & \(S_1 + 2J_{a,1} + J_{b,1}\) & $\Ga^6\rtimes (\Gm\times (\Sp_4(k)\times \SL_2(k)))$ & $T$ \\
16 & \(S_1 + J_{a,1} + J_{b,1} + J_{c,1}\) & $\Ga^6\rtimes (\Gm\times (\SL_2(k)\times \SL_2(k)\times \SL_2(k)))$ & $S_3$ preimage \\
\end{longtable}

\endgroup

\subsection{How the formalization is used}
\label{subsec:formalization-protocol}

For each table row there is a public Lean namespace
\path|Wedge2Formalization.Paper.Nn.Rowm| and a matching source file
\path|Wedge2Formalization/Paper/Nn/Rowm.lean| in the companion repository.
This is the referee-facing layer of the formalization; the lower-level
coordinate files are not needed for a row-by-row check.

Inside that namespace, the declarations \path|paperRep1| and
\path|paperRep2| are the literal basis vectors printed in the geometric orbit
table. When the proof uses a more convenient internal working model, the
lemmas \path|paperRep1_transport| and \path|paperRep2_transport| identify the
displayed representative with that internal model.

The theorem \path|mem_K_table_iff| is the public kernel statement for the row,
and the theorem \path|quotient_image| identifies the image \(Q_L\) of the
coefficient-side action. Together with the representative transport lemmas,
these declarations recover the stabilizer line printed in the tables. Appendix
~\ref{app:lean-crosswalk} records the row-to-namespace correspondence.
